\documentclass[sigconf, nonacm]{acmart}


\usepackage{draftwatermark}
\SetWatermarkText{AI-Generated}
\SetWatermarkScale{0.5}
\SetWatermarkColor[gray]{0.6}
\AtBeginDocument{\providecommand\BibTeX{{Bib\TeX}}}

\copyrightyear{2025}
\acmYear{2025}
\setcopyright{none}
\acmConference[arXiv preprint]{arXiv preprint}{2025}{}
\acmBooktitle{}
\acmDOI{}
\acmISBN{}
\settopmatter{printacmref=false}         
\renewcommand\footnotetextcopyrightpermission[1]{} 

\usepackage{multirow}
\usepackage{enumitem}
\usepackage[skip=2pt]{caption}
\usepackage{bm}
\graphicspath{{figures/}}


\begin{document}
\title{FoodLogAthl-218: Constructing a Real-World Food Image Dataset\\ Using Dietary Management Applications}

\author{Anonymous Author}

\begin{abstract}
Food recognition models for dietary logging are usually trained and benchmarked on web-sourced photographs that are attractive, class-balanced and shot one dish at a time, unlike the photographs people record in a meal-logging app. We introduce FoodLogAthl-218, a benchmark built bottom-up from user-submitted meal records of the FoodLog Athl application: labels are assigned after the photos are collected and then cleaned and consolidated, preserving natural class imbalance, casual imaging, multi-dish meals and temporal continuity. CLIP-based outlier filtering, manual all-or-none review and GPT-4o-assisted label consolidation reduce 11,460 raw images from 632 users to 6,925 images with 14,349 dish bounding boxes from 513 users in 218 classes. Three tasks are defined: closed-set dish classification, continual user-aware adaptation along the time axis, and menu-aware classification in multi-dish images. Ten multimodal models and CLIP-style encoders reach at most 0.424 Recall@1 (GPT-4o) and 0.600 Recall@5 (CLIP ViT-L/14), and the ranking depends on the metric. Feature-space statistics indicate larger intra-class spread and more class overlap than UEC-Food 256. Fine-tuning on accumulated logs raises weekly accuracy from roughly 0.25 to 0.5 while zero-shot accuracy does not improve, a user-specific adapter gains about 50 points for one heavy logger, and adding whole-image description text trades recall for precision (F1 0.159 to 0.230, Recall@1 0.366 to 0.324). The benchmark covers one application and reports no inter-annotator agreement; we state what it can and cannot support.
\end{abstract}

\maketitle

\section{Introduction}
Image-based dietary logging has users photograph what they eat and the system recognise the dishes, so that intake can be tracked~\cite{aizawa2015foodlog,meyers2015im2calories,nakamoto2022foodlog}. Progress on the underlying recognition problem has mostly been measured on datasets collected from the web for a predefined list of categories, such as Food-101~\cite{bossard2014food}, UEC-Food 256~\cite{kawano2015automatic}, ISIA Food-500~\cite{min2020isia} and Food2K~\cite{min2023large}. Their construction shapes what they measure: images are searched for a given label and tend to be photographed to be shared, so they show one attractive dish, are often class-balanced, are independent of one another, and carry no information about who ate what and when. A model that scores well on them has not necessarily been tested under logging conditions.

We address this gap with FoodLogAthl-218, constructed from records of the FoodLog Athl application~\cite{nakamoto2022foodlog}. The key design decision is the order of construction (Figure~\ref{fig:distribution}): rather than choosing categories and collecting images for them, we start from real user logs, labeled by the user at logging time with help from the application, and clean and consolidate the labels afterwards. This keeps what web collection removes: a heavy-tailed class distribution, casual photographs, several dishes per photograph, and a personal meal history with time stamps and anonymized user identifiers. The benchmark measures how well recognisers cope with these conditions in three settings (Figure~\ref{fig:tasks}); it does not measure nutrient or portion estimation, or recognition of cuisines outside those in the application.

Evaluating ten representative systems gives four findings. (i) Closed-set classification over 218 classes is far from solved: no system exceeds 0.424 Recall@1, and the ranking changes with the metric, as contrastive encoders lead in Recall@5 while the strongest LMM leads in Recall@1 and F1. (ii) In feature space the benchmark is more spread out and less separable than UEC-Food 256. (iii) Fine-tuning on accumulated logs improves accuracy over time while the zero-shot model does not, and a per-user adapter helps a heavy logger substantially. (iv) Whole-image description text changes the error profile rather than uniformly improving it: precision and F1 rise while recall falls. Our contributions are:
\begin{itemize}[leftmargin=1.2em,itemsep=0pt,topsep=2pt]
  \item FoodLogAthl-218: 6,925 images, 14,349 dish bounding boxes and 218 consolidated classes from real dietary-app logs, with natural imbalance and time and user metadata, collected with ethics approval and informed consent.
  \item A reproducible mixed automated--human curation pipeline (CLIP outlier filtering, manual all-or-none review, GPT-4o label grouping with human adjudication).
  \item Three tasks with a common protocol (Top-1/3/5 recall and F1): closed-set, continual user-aware, and menu-aware classification.
  \item Baselines with CLIP-style encoders and LMMs, a diversity comparison with UEC-Food 256, and an analysis of what success and failure on the benchmark mean.
\end{itemize}

\begin{figure}[t]
  \centering
  \includegraphics[width=0.78\columnwidth]{ACM_task.jpg}
  \caption{The three evaluation tasks. Task~1: closed-set classification of a dish crop; a free-form description from a large multimodal model (LMM) is embedded and matched to class-name embeddings. Task~2: continual user-aware adaptation, in which a model is repeatedly updated with accumulated logs and tested on later data. Task~3: menu-aware classification, in which a description of the whole meal photograph is given as context for one dish.}
  \label{fig:tasks}
\end{figure}

\begin{figure*}[t]
  \centering
  \includegraphics[width=0.88\textwidth]{distribution_5.jpg}
  \caption{Traditional web-collected food image datasets (top) versus FoodLogAthl-218 (bottom): data source, label-after-image collection, natural class imbalance, low aesthetic bias, time-series continuity and multi-dish meals. The bar charts are schematic, not measured data.}
  \label{fig:distribution}
\end{figure*}

\section{Related Works}
\paragraph{Food image datasets.}
Early food recognition datasets were small, for example PFID~\cite{chen2009pfid} and Japanese-food datasets of up to 256 categories~\cite{hoashi2010image,joutou2009food,kawano2015automatic}. Web-scale resources followed: Food-101~\cite{bossard2014food} and UPMC Food-101~\cite{wang2015recipe}, ChineseFoodNet~\cite{chen2017chinesefoodnet}, VegFru~\cite{hou2017vegfru}, FoodAI-756~\cite{sahoo2019foodai}, ISIA Food-500~\cite{min2020isia}, Food-500 Cap~\cite{ma2023food} and Food2K with 2,000 classes and over one million images~\cite{min2023large}. Recipe1M+~\cite{marin2019learning}, Nutrition5k~\cite{thames2021nutrition5k} and the Cookpad image dataset~\cite{harashima2017cookpad} pair images with recipes, nutrition or recipe-sharing content, and FoodX-251~\cite{kaur2019foodx}, a benchmark on the representation of food images~\cite{farinella2014benchmark} and a food segmentation benchmark~\cite{wu2021large} address fine-grained classification, representation and segmentation. Apart from datasets built from experiment participants such as UNICT-FD1200~\cite{farinella2016retrieval}, the classification datasets of Table~\ref{Food image dataset table} obtain images from the web or from controlled collections, and none is derived from the logs that users record for dietary management. FoodLogAthl-218 complements them with such records, labeled after collection.

\paragraph{Food logging and personalization.}
Multimedia food-logging systems~\cite{aizawa2015foodlog,nakamoto2022foodlog}, including recognition engines built for dietary logging~\cite{merler2016snap}, and commercial applications~\cite{asken,myfitnesspal} accumulate personal meal histories. Earlier work on such logs exploited personal dietary tendencies~\cite{Aizawa2013FoodBalance}, extracted food categories from meal names~\cite{Amano2015FoodCategory}, and built personalized classifiers~\cite{flg469,Horiguchi2018Personalized}, including by sharing records between users~\cite{Kim2021Boosting}. To our knowledge, none was accompanied by a public benchmark with time ordering and user identity; Task~2 makes this setting an explicit evaluation target, measuring accuracy along the time axis as logs accumulate rather than on a static split.

\paragraph{Vision-language models for food.}
Contrastive encoders such as CLIP~\cite{CLIP} and EVA-CLIP~\cite{sun2023eva} allow zero-shot classification, and large multimodal models have been adapted to food tasks including recipe generation and food assistance~\cite{yin2023foodlmm,jiao2024rode,Imajuku2025FoodMLLM}. WorldCuisines benchmarks multilingual and multicultural food question answering~\cite{winata2024worldcuisines}. We evaluate LMMs through lmms-eval~\cite{lmms_eval2024,zhang2024lmmsevalrealitycheckevaluation} and adapt them with LoRA~\cite{hu2022lora}. FoodLogAthl-218 instead targets recognition of individual dishes in photographs taken for personal logging, with time and user structure.

\section{FoodLogAthl-218 Dataset}
\subsection{Source and design principles}
FoodLog Athl is a dietary management application in which users photograph meals and the system proposes dish regions and names~\cite{nakamoto2022foodlog}. Each record has an anonymized user identifier, date and time, dish names chosen from a database of about 140K items, portion size and the initial detections; users could correct labels. We use records from May 2023 to October 2024: initially 11,460 images with 61,414 bounding boxes and 5,889 raw labels from 632 users. Four principles guide the design: \emph{realism} (actual logs, no web search for predefined classes); \emph{natural imbalance} (frequencies are not equalised); \emph{meal structure} (the image is the atomic unit, multi-dish photographs are kept whole); and \emph{temporal and user grounding} (time stamps and anonymous user identifiers are kept).

\subsection{Curation pipeline}
\input{tables/table_zn.tex}
Table~\ref{data cleaning steps table} lists the steps. (1--2)~Records without an image or a bounding box are removed, leaving 11,285 images and 26,065 boxes in 3,071 classes. (3)~Classes with fewer than 10 samples are dropped (9,121 images, 21,944 boxes, 1,180 classes). (4)~Features from CLIP ViT-B/16~\cite{CLIP} flag samples more than 1.5 standard deviations from their class centroid, which are removed (7,537 images, 17,513 boxes). (5)~Reviewers manually fix or remove invalid boxes and labels; since a meal is the unit of the benchmark, an image with an invalid sample is removed entirely (all-or-none) so that remaining multi-dish photographs stay complete, leaving 6,925 images and 14,349 boxes from 513 users. (6)~GPT-4o groups synonymous labels and manual refinement of the groups merges the 1,180 classes into 218. Figure~\ref{fig:errors} shows the defects the review targets: incorrect labels, incorrect boxes and inappropriate images. Personal information is removed in two manual stages; data collection had ethics approval and informed consent.

\begin{figure}[t]
  \centering
  \includegraphics[width=\columnwidth]{error_example.jpg}
  \caption{Examples of erroneous data in the raw logs: incorrect label, incorrect bounding box and inappropriate images.}
  \label{fig:errors}
\end{figure}

\subsection{Characteristics of the resulting data}
Figure~\ref{fig:thumbs} shows typical samples: lunch boxes, packaged food, partly eaten dishes, clutter, and several dishes per photograph (14,349 boxes in 6,925 images, about 2.1 per image). Figure~\ref{fig:hist} shows the images per class on a logarithmic axis. No class had fewer than 10 samples before consolidation; class sizes range from 10 to roughly 2,000 images (read from the figure), with 16 classes at the 10-image end and a long tail of frequent classes, so the distribution is strongly imbalanced. Table~\ref{Food image dataset table} places the dataset among existing ones. Its sample count is in dish bounding boxes (crops from 6,925 images) whereas the others count images, and its class count is small because labels were merged to remove synonyms.

\begin{figure}[t]
  \centering
  \includegraphics[width=0.7\columnwidth]{thumbnail_re.jpg}
  \caption{Examples of data in FoodLogAthl-218: casual photographs with packaged food, lunch boxes and multiple dishes.}
  \label{fig:thumbs}
\end{figure}

\begin{figure}[t]
  \centering
  \includegraphics[width=0.88\columnwidth]{histogram_linear_pdf.jpg}
  \caption{Number of dish types (vertical, linear) with a given number of images per class (horizontal, logarithmic); the arrows mark classes of different sizes (unlabeled in the source figure). The x-axis counts images, and no class has more than about 2,000.}
  \label{fig:hist}
\end{figure}

\input{tables/table_os.tex}

\subsection{Validation of the benchmark}
\paragraph{Diversity compared with a web-collected benchmark.}
We compare the benchmark with UEC-Food 256~\cite{kawano2015automatic}, which has a similar class count, using features from CLIP~\cite{CLIP}, EVA-CLIP~\cite{sun2023eva} and Swin Transformer~\cite{liu2021swin}: the trace of the feature covariance as a measure of spread, and the Silhouette coefficient~\cite{rousseeuw1987silhouettes}, Davies--Bouldin index~\cite{davies1979cluster} and Calinski--Harabasz index~\cite{calinski1974dendrite} of the class partition as measures of separation. Table~\ref{table:diversity_metrics_comparison} shows that for all three extractors FoodLogAthl-218 has a larger trace (CLIP 81.70 versus 71.21), a lower Silhouette coefficient (CLIP $-0.0210$ versus $0.0335$; negative for CLIP and Swin features), a higher Davies--Bouldin index and a lower Calinski--Harabasz index. All four indicators agree in direction for all three extractors, suggesting larger intra-class spread and more overlap between classes. These are descriptive statistics from one run without confidence intervals, and the datasets differ in class count (218 versus 256), sample count and taxonomy, so the comparison supports the statement that classes are less separable in these feature spaces, not a calibrated difficulty ranking.

\input{tables/table_xj.tex}

\paragraph{Label and annotation quality.}
Quality control combines automated outlier detection, manual correction of labels and boxes, and human adjudication of the GPT-4o synonym groups. We do \emph{not} report inter-annotator agreement, so residual label noise is unquantified; labels originate from user choices and consolidation, which also makes some boundaries between merged classes subjective. The CLIP-based filter removed samples far from class centroids and may preferentially remove atypical or hard examples, so Task~1 difficulty may be underestimated relative to the unfiltered logs; we have not measured this.

\section{Benchmark Tasks and Experimental Setup}
All tasks use the dish bounding boxes and the 218 classes. Performance is reported as Top-$K$ recall (R@$K$, $K=1,3,5$: the fraction of samples whose true class is among the $K$ highest-ranked labels) and, where stated, F1 and Precision@1 (P@1). Data splits, prompts, LoRA settings, test-set sizes and the F1 averaging are not specified in the material available to us, and we do not state them.

\paragraph{Task 1: closed-set dish classification.}
The input is the crop of one dish and the output a ranking over the 218 classes. CLIP ViT-L/14~\cite{CLIP} and EVA-CLIP ViT-g~\cite{sun2023eva} rank class names by image--text similarity. LMMs generate text, so we prompt for a free-form description of the dish, embed it with intfloat/e5-large-v2~\cite{wang2023e5}, and rank labels by cosine similarity between their embeddings and the description. We run models through lmms-eval~\cite{lmms_eval2024} and evaluate GPT-4o~\cite{achiam2023gpt}, Qwen2.5-VL-7B-Instruct, Qwen2-VL-7B-Instruct, Qwen2-VL-3B-Instruct, and the LLaVA-family models llava-onevision-qwen2-7b-ov, llava-v1.6-34b, llava-v1.5-7b and llava-v1.6-mistral-7b~\cite{liu2024improved,liu2024visual}.

\paragraph{Task 2: continual user-aware adaptation.}
Logs arrive over time and a deployed recogniser can be updated with the logs accumulated so far. Task~2 orders the data by time, trains on the logs up to a point and evaluates on the following data (Figure~\ref{fig:tasks}). We report accuracy per week together with a weighted one-month moving average ($\pm2$ weeks) whose weights are proportional to the weekly data volume. We compare the zero-shot model with full fine-tuning and LoRA fine-tuning~\cite{hu2022lora} on all users' data (only the all-users LoRA model is stated to be Qwen2.5-VL-3B; the base models of the full fine-tuning curve and of the User~A adapter are not stated in the material). A user-specific setting trains LoRA only on the logs of one heavy logger (``User~A'') and compares it with zero-shot and all-users LoRA.

\paragraph{Task 3: menu-aware classification.}
The model first produces a rich, objective description of the whole photograph (the ``menu'') and then classifies the cropped dish given that description (Figure~\ref{fig:desc}). We compare a 7B Qwen-VL model (labeled Qwen2.5-VL-7B-Instruct in the experiment table) with and without the description, reporting F1, P@1 and R@1/3/5 and the number of predictions that change from wrong to right and vice versa.

\section{Benchmark Results and Analysis}
\subsection{Task 1: closed-set classification}
\input{tables/table_kf.tex}
Table~\ref{table:classification_comparison} shows GPT-4o best in F1 (0.3732), R@1 (0.4243) and R@3 (0.4991) and CLIP ViT-L/14 best in R@5 (0.6000). No system identifies even half of the dishes at rank 1, and 40\% are outside the top five of CLIP ViT-L/14, so the benchmark is not saturated. The ranking depends on the metric: CLIP ViT-L/14 and EVA-CLIP reach R@1 of only 0.2567 and 0.2654, below GPT-4o and the 7B Qwen and LLaVA-OneVision models (0.363 to 0.424), yet their R@5 (0.6000 and 0.5944) is above all LMMs, including GPT-4o (0.5221). An LMM must name a dish in free text that is then matched to labels, whereas the encoders compare the crop with all 218 class names directly, so we do not attribute the gap to a single cause. Qwen2-VL-3B is clearly below its 7B counterpart (R@1 0.2663 versus 0.3656), the older LLaVA models are weakest (R@1 0.1932 to 0.2432), and F1 is far below R@1 for most LMMs, for example 0.1594 versus 0.3656 for Qwen2-VL-7B-Instruct (the gap is smallest for GPT-4o, 0.3732 versus 0.4243), whereas the F1 of the encoders (about 0.24) is close to their R@1 (about 0.26). Each model was evaluated once, so we report no variance.

\subsection{Task 2: continual user-aware adaptation}
\begin{figure}[t]
  \centering
  \includegraphics[width=0.9\columnwidth]{all_data_1_35_with_legend.jpg}
  \caption{Incremental accuracy of zero-shot, full fine-tuning and LoRA fine-tuning on all users. Dashed lines: weekly values; solid lines: weighted moving average.}
  \label{fig:all}
\end{figure}
Figure~\ref{fig:all} shows weekly accuracy for all users from June 2023 to early January 2024 (the figures of Task~2 cover only part of the May 2023 to October 2024 span of the logs); values below are read from the plot and are approximate. The moving average of the zero-shot model stays near 0.23 until August 2023 and then drifts down to about 0.16, while both fine-tuned models rise from about 0.24 to about 0.5 by January 2024 (full fine-tuning about 0.51, LoRA about 0.49), a gap of over 30 percentage points to zero-shot. Full fine-tuning is above LoRA for most of the period and the two are close at the end. Weekly values fluctuate widely (roughly 0.03 to 0.6) and no confidence intervals are available, so the moving average is the more reliable summary. We have not analysed why the zero-shot curve declines (a shift in dish mix or in the active users is possible).

\begin{figure}[t]
  \centering
  \includegraphics[width=0.9\columnwidth]{user_zero_lora_ma_wide_legend.jpg}
  \caption{Accuracy over time for User~A: zero-shot, LoRA fine-tuned on all users, and LoRA fine-tuned on User~A only. X ticks are every four weeks and labeled by month only, so 2023/10 appears twice. Dashed: weekly values; solid: weighted moving average.}
  \label{fig:user}
\end{figure}
Figure~\ref{fig:user} turns to a single heavy logger (its tick label ``2023/10'' appears twice on the axis, so month readings are approximate). Zero-shot accuracy for User~A stays at about 0.07 to 0.14. The user-specific LoRA model starts at about 0.17, below the all-users LoRA model (about 0.30 in the first weeks), overtakes it around the middle of the period and stays ahead for several months at roughly 0.6 versus 0.4 to 0.55, about 50 percentage points above zero-shot. Near the end the all-users model catches up and is slightly higher (about 0.65 versus 0.57). User-specific adaptation thus pays off once enough personal data have accumulated, but its advantage over all-users adaptation is temporary in this example. With one user and no repeated runs, this is an existence demonstration of personalization gains, not an estimate of their typical size.

\subsection{Task 3: menu-aware classification}
\begin{figure}[t]
  \centering
  \includegraphics[width=0.82\columnwidth]{description_eg.jpg}
  \caption{Description-guided classification. Left: the whole-meal description corrects an error on a dumpling crop. Right: a description that suggests a drink induces an error on a correctly classified crop.}
  \label{fig:desc}
\end{figure}
\input{tables/table_vs.tex}
Table~\ref{table:desc_guided} shows that the description raises F1 from 0.1594 to 0.2304 and P@1 from 0.1019 to 0.1788, while R@1, R@3 and R@5 fall (R@1 from 0.3656 to 0.3240). The study reports that the description corrected 687 errors but caused 1,222 new ones. Figure~\ref{fig:desc} shows both directions: in a success the description mentions gyoza and a pan-fried appearance and the crop is classified as dumplings instead of okonomiyaki; in a failure the description speculates that a cup holds juice or a smoothie and a crop correctly classified as vegetable juice without context becomes egg soup. Menu context thus changes the precision--recall profile and is not a uniform gain. A consistency note: the ``w/o description'' row of Table~\ref{table:desc_guided} coincides numerically with the Qwen2-VL-7B-Instruct row of Table~\ref{table:classification_comparison} (F1 0.1594; R@1/3/5 0.3656/0.4590/0.4781), whereas the Qwen2.5-VL-7B-Instruct row there differs (F1 0.2349; R@1/3/5 0.3684/0.4609/0.4859). We report the Task~3 values exactly as supplied, cannot resolve from the material which label is correct, and rely only on the comparison between the two rows.

\section{Discussion}
\paragraph{What the benchmark offers.}
Task~1 shows that the strongest evaluated systems leave a large share of dishes misidentified, Task~2 that adapting to accumulated logs yields more than 30 percentage points over the zero-shot model in this setting, and Task~3 that context can trade recall for precision. We recommend reporting both R@1 and R@5 for Task~1, since they rank systems differently.

\paragraph{What cannot be concluded.}
The data come from one application and one user population; the examples are predominantly Japanese-style meals and we report no cuisine breakdown, so results should not be read as covering all cuisines or logging behaviors, although Table~\ref{Food image dataset table} lists the category as international. The 218 classes result from GPT-4o-assisted and manual merging of user-chosen labels, so class boundaries reflect those decisions and some classes may be ambiguous. Inter-annotator agreement is not reported. The CLIP-based filter and all-or-none removal discard difficult or atypical meals, so the benchmark does not represent all logged meals, and CLIP-family models may be advantaged by the use of CLIP features in curation, which we have not tested. The diversity statistics compare against a single dataset without uncertainty estimates. Tasks~2 and~3 rest on one model each and, for personalization, one heavy user, and the Task~3 comparison cannot separate the value of context from that of an extra generation step. Use of the personal meal records should respect the consent and ethics conditions under which they were obtained. A high score would show competence at dish classification on logged meals, not general food understanding or nutritional accuracy.

\section{Conclusion}
FoodLogAthl-218 is a benchmark of 6,925 images, 14,349 dish bounding boxes and 218 classes built from real dietary-app records by labeling after collection, with tasks for closed-set classification, continual user-aware adaptation and menu-aware classification. Current LMMs and CLIP-style encoders reach at most 0.4243 R@1 and 0.6000 R@5; the benchmark is less separable than UEC-Food 256 in feature space; fine-tuning on accumulated logs lifts accuracy from about 0.25 to about 0.5; a user-specific adapter helps a heavy logger by about 50 points; and menu description trades recall for precision. The single-source origin, missing agreement statistics and unquantified filtering effects bound what the benchmark can claim.

\bibliographystyle{ACM-Reference-Format}

\bibliography{references}

\end{document}
